% !TeX root = ../../../main.tex
\subsection{\bilingual{First-Moment Reduction for the Speed Bound}{速率界的一阶矩约化}}
\label{section:speed-first-moment-reduction}

\bilingual{
    We isolate the step needed to replace the fourth-moment argument in the
    proof of \Cref{theorem:speed}.  A first moment alone does not give a rate
    for the integral over an exceptional event.  The required rate follows
    when the continuation used on that event is sampled from a reserve subtree
    independent of the information that declares the event.
}{
    下面分离出在\Cref{theorem:speed} 的证明中替代四阶矩论证所需的步骤. 单独的一阶矩不能给出坏事件上积分的收敛速率. 若在该事件上使用的延伸取自一个独立于事件判定信息的备用子树, 则可得到所需速率.
}

\begin{lemma}[\bilingual{Independent exceptional-event estimate}{独立坏事件估计}]
    \label{lemma:independent-exceptional-L1}
    \bilingual{
        Let $Y$ be integrable and let $A$ be an event independent of $Y$. Then
    }{
        设 $Y$ 可积, 且事件 $A$ 与 $Y$ 独立. 则
    }
    \begin{equation*}
        \mathbb E\!\left[|Y|\mathbf 1_A\right]
        =\mathbb E[|Y|]\mathbb P(A).
    \end{equation*}
    \bilingual{
        Consequently, if $Y_1,\ldots,Y_n$ are integrable and each is independent
        of $A$, then
    }{
        因而, 若 $Y_1,\ldots,Y_n$ 均可积且分别与 $A$ 独立, 则
    }
    \begin{equation*}
        \mathbb E\!\left[
          \left(Cn+\sum_{j=1}^n|Y_j|\right)\mathbf 1_A
        \right]
        =\left(Cn+\sum_{j=1}^n\mathbb E|Y_j|\right)\mathbb P(A).
    \end{equation*}
\end{lemma}

\begin{proof}
    \bilingual{
        Apply the defining integral factorization for independent random
        variables to $|Y|$ and $\mathbf 1_A$.  The finite-sum identity follows
        by linearity of expectation.  Notice that mutual independence of the
        $Y_j$ is unnecessary for this estimate.
    }{
        对独立随机变量 $|Y|$ 与 $\mathbf 1_A$ 应用积分乘积分的分解公式. 有限和恒等式由期望的线性性得到. 此估计并不要求各 $Y_j$ 相互独立.
    }
\end{proof}

\bilingual{
    In the restart construction, let $E_N$ be the successful-coupling event
    and let $n=\ell_N=O(\log^3N)$.  If the continuation on $E_N^{\mathrm c}$
    is read from a pre-sampled reserve lineage disjoint from the coordinates
    inspected in declaring $E_N$, the branching property makes its increments
    independent of $E_N$.  With $\mathbb P(E_N^{\mathrm c})\le N^{-\varepsilon}$
    and $\mathbb E[(\Xi_1)_+]<\infty$, the preceding lemma gives
}{
    在重启构造中, 令 $E_N$ 为成功耦合事件并令
    $n=\ell_N=O(\log^3N)$. 若在 $E_N^{\mathrm c}$ 上的延伸读取自一条预采样备用谱系, 且该谱系与判定 $E_N$ 时已检查的坐标不交, 则分支性质说明其增量与 $E_N$ 独立. 利用
    $\mathbb P(E_N^{\mathrm c})\le N^{-\varepsilon}$ 与
    $\mathbb E[(\Xi_1)_+]<\infty$, 上述引理给出
}
\begin{equation*}
    \mathbb E\!\left[
      \left(C\ell_N+\sum_{j=1}^{\ell_N}(R_j)_+\right)
      \mathbf 1_{E_N^{\mathrm c}}
    \right]
    =O(\log^3N)N^{-\varepsilon}
    =o(\log N).
\end{equation*}
\bilingual{
    This is the $L^1$ estimate that would replace the Cauchy--Schwarz bound in
    the upper speed proof once the required reserve continuation is
    constructed. The argument applies to the restart process only after one
    proves that its actual exceptional event is independent of the reserve
    continuation and that the latter bounds the failed continuation
    pathwise. The quantitative zero-start tube estimate, including the
    Mogul'ski\u i asymptotic, and the final speed-limit assembly also remain
    open.
}{
    备用延伸构造完成后, 这条一阶矩估计可替代速率上界中的 Cauchy--Schwarz 界. 要用于本文的重启过程, 还需证明实际坏事件与备用延伸独立, 并证明备用路径逐点控制失败路径. 零起点管道的小偏差估计及最终速率拼接也仍待完成.
}
